% Figure: radial and hoop stress in a solid rotating disk of uniform
% thickness, normalised by rho omega^2 a^2 and plotted against r/a.
% Closed form (plane stress, free rim):
%   sigma_r/(rho omega^2 a^2)     = (3+nu)/8 * (1 - (r/a)^2)
%   sigma_theta/(rho omega^2 a^2) = (3+nu)/8 - (1+3 nu)/8 * (r/a)^2
% Both peak at the centre where they are equal to (3+nu)/8; the hoop stays
% finite and largest at the axis. Uses nu=0.3.
\begin{figure}[htbp]
\centering
\begin{tikzpicture}
\begin{axis}[
  width=11cm, height=6.5cm,
  xlabel={normalised radius $\ r/a$},
  ylabel={stress / $(\rho\omega^2 a^2)$},
  xmin=0, xmax=1.0, ymin=0, ymax=0.45,
  axis lines=left,
  grid=major, grid style={enggray!20},
  tick label style={font=\scriptsize},
  label style={font=\small},
  legend style={font=\scriptsize, at={(0.98,0.98)}, anchor=north east, draw=enggray!40},
]
% radial stress: (3.3/8)(1 - x^2), nu=0.3 -> (3+nu)/8 = 0.4125
\addplot[engblue, very thick, domain=0:1, samples=90]
  {0.4125*(1 - x*x)};
\addlegendentry{radial $\sigma_r$}
% hoop stress: 0.4125 - (1+0.9)/8 x^2 = 0.4125 - 0.2375 x^2
\addplot[engred, very thick, domain=0:1, samples=90]
  {0.4125 - 0.2375*x*x};
\addlegendentry{hoop $\sigma_\theta$}
% mark the common centre value
\addplot[engamber, only marks, mark=*, mark size=1.6pt] coordinates {(0,0.4125)};
\node[engamber, font=\scriptsize, anchor=west] at (axis cs:0.04,0.4125)
  {centre $\tfrac{3+\nu}{8}$};
% mark the rim hoop value (0.175 at x=1)
\addplot[engred, only marks, mark=*, mark size=1.4pt] coordinates {(1.0,0.175)};
\node[engred, font=\scriptsize, anchor=south east] at (axis cs:0.98,0.185)
  {rim hoop};
\end{axis}
\end{tikzpicture}
\caption[Stress in a spinning disk]{Radial and hoop stress in a solid disk
of uniform thickness spinning at angular speed $\omega$, from the plane-stress
solution with a free rim, normalised by $\rho\omega^2 a^2$ and plotted against
the normalised radius ($\nu=0.3$). The centrifugal body force $\rho\omega^2 r$
is a load that grows with radius, yet both stresses are largest at the centre,
where they meet at the common value $(3+\nu)/8$ (amber). The radial stress
(blue) falls to zero at the free rim; the hoop stress (red) stays finite and
carries the largest tension at the axis, which is why a bursting flywheel or
turbine disk fails from the centre outward and why a central bore, which
removes the very material that carries the peak, nearly doubles the bore hoop
stress. The whole field is set by $\rho\omega^2 a^2$ and the geometry, an
inertial load with no external force applied at any boundary.}
\label{fig:vol03ch09:spinning-disk}
\end{figure}
